\chapter{Supersymmetry as Odd Symmetry and Superspace}

Supersymmetry enters this lecture with an explicit warning. It is not presented as a tidy embellishment of familiar symmetry, but as one of the most abstract constructions in physics. Susskind therefore refuses to begin with the final algebra. He first backs up and reviews the odd variables out of which the subject is built. These notes follow that order on purpose. We first learn how Grassmann quantities behave, then recall the ordinary language of continuous symmetries, then build the simplest supersymmetry algebra, and only afterward turn to the geometric interpretation in terms of superspace. The key surprise is that the odd generators do not close on another internal generator. Their anticommutator closes on the Hamiltonian itself.

\section{Why Supersymmetry Begins with Grassmann Numbers}

Supersymmetry is described as a symmetry of space-time, but of a very bizarre kind. Rotations, translations, and the ordinary internal symmetries of particle physics are familiar enough: they act either on ordinary coordinates or on particle labels. Supersymmetry does something stranger. It mixes bosonic and fermionic states, and in order to write it naturally one is driven to odd variables that are not ordinary numbers at all.

That is the reason for the apparently backward beginning. Before saying what the symmetry is, the lecture teaches us the language in which it will be expressed. Grassmann variables are not introduced as decorative formalism. They are the minimal machinery needed for the later algebra. Without them, the odd generators \(Q\) and \(\bar Q\) would look unmotivated; with them, the later manipulations become a continuation of rules we have already learned.

\section{Grassmann Numbers, Functions, and Calculus}

\paragraph{Odd and even elements.}
Let \(\epsilon_i\) denote Grassmann variables. They are odd elements of the algebra, and they anticommute:
\begin{equation}
\epsilon_i\epsilon_j+\epsilon_j\epsilon_i=0.
\end{equation}
In particular,
\begin{equation}
\epsilon_i^2=0.
\end{equation}
Ordinary real or complex numbers are even elements. The parity rules are
\begin{equation}
\text{odd}\cdot\text{odd}=\text{even},
\qquad
\text{odd}\cdot\text{even}=\text{odd},
\qquad
\text{even}\cdot\text{even}=\text{even}.
\end{equation}
Even elements commute with everything. The only genuine anticommutation comes from moving one odd element past another.

This much already explains why the subject looks unusual. The square of an odd Grassmann variable vanishes, but the variable itself need not vanish. That single fact drastically truncates the algebra.

\paragraph{Functions of Grassmann variables.}
A function of one Grassmann variable cannot be an infinite power series in the usual sense, because the square of the variable is zero. Thus
\begin{equation}
f(\epsilon)=a+b\,\epsilon.
\end{equation}
With two independent odd variables the most general form is
\begin{equation}
f(\epsilon_1,\epsilon_2)=a+b_1\epsilon_1+b_2\epsilon_2+c\,\epsilon_1\epsilon_2.
\end{equation}
Nothing higher survives. If we pass to a complex pair \(\theta,\bar\theta\), the same structure becomes
\begin{equation}
f(\theta,\bar\theta)=a+\bar b\,\theta+\bar\theta\,b+c\,\bar\theta\theta,
\end{equation}
up to the usual sign conventions associated with ordering the odd variables.

The lecture makes a practical point here. There are not very many functions of finitely many Grassmann variables. One is working with a small polynomial algebra, not with a large analytic zoo. That is why Susskind keeps describing the formalism as bookkeeping. It is a strange bookkeeping, but still bookkeeping.

\paragraph{Even and odd functions.}
A function built from Grassmann variables is usually taken to have definite parity. If \(f\) is even, then the coefficients must alternate in parity so that each term is even. If \(f\) is odd, the pattern reverses. One can invent mixed expressions, but in the lecture the natural objects are always taken to be either even or odd.

That point will matter later. Ordinary observables are even. Supersymmetry generators will be odd.

\paragraph{Differentiation and sign bookkeeping.}
Let \(\theta\) be a Grassmann variable. Differentiation with respect to \(\theta\) is simple once the ordering is handled correctly. If a term does not depend on \(\theta\), then
\begin{equation}
\partial_\theta a=0.
\end{equation}
If the \(\theta\)-dependence is written in standard position, then
\begin{equation}
\partial_\theta(a+\theta b)=b.
\end{equation}
But one must be careful when odd coefficients sit on the other side. If \(b\) is odd, then
\begin{equation}
b\theta=-\theta b,
\end{equation}
so before differentiating we bring \(\theta\) next to the derivative and keep the sign:
\begin{equation}
\partial_\theta(b\theta)=-b.
\end{equation}

This is exactly the kind of local obstacle the lecture wants us to feel. The calculus is not difficult in itself; the difficulty is remembering that every time one odd element passes another, the sign changes.

For two odd variables \(\theta\) and \(\bar\theta\), the same issue appears in a slightly richer example. If
\begin{equation}
f(\theta,\bar\theta)=a+\bar b\,\theta+\bar\theta\,b+c\,\bar\theta\theta,
\end{equation}
then differentiating with respect to \(\theta\) requires us to move each \(\theta\) into the correct place before acting. In an even example, this produces the sign pattern
\begin{equation}
\partial_\theta f=-\,\bar b-c\,\bar\theta,
\end{equation}
while in an odd example the signs come out differently. The lecture's lesson is not the specific formula but the method: reorder first, then differentiate.

\paragraph{A bridge to operator algebra.}
Susskind now compares Grassmann calculus with ordinary calculus. For an ordinary coordinate \(x\),
\begin{equation}
[\partial_x,x]=1.
\end{equation}
The corresponding Grassmann identity is obtained by acting on a test function
\begin{equation}
f(\theta)=a+\theta b.
\end{equation}
We compute
\begin{align}
\theta\,\partial_\theta f &= \theta b, \\
\partial_\theta(\theta f) &= \partial_\theta(\theta a+\theta^2 b)=a,
\end{align}
because \(\theta^2=0\). Adding the two orders gives
\begin{equation}
\bigl(\partial_\theta\theta+\theta\partial_\theta\bigr)f
=
a+\theta b
=
f.
\end{equation}
Since this works for an arbitrary function of one Grassmann variable,
\begin{equation}
\{\partial_\theta,\theta\}=1.
\end{equation}

This identity is one of the lecture's turning points. We are no longer merely listing Grassmann rules. We have reached an operator relation of the same flavor as the commutators that organize symmetry generators in ordinary quantum mechanics.

\paragraph{Integration.}
The integral calculus is only stated, not derived in detail, but the statement is striking:
\begin{equation}
\int d\theta\,f(\theta)=\partial_\theta f(\theta).
\end{equation}
The lecture motivates this by saying that if one insists on integration by parts, the Grassmann integral must coincide with the Grassmann derivative. Once again the point is that the whole calculus is tiny compared with the ordinary one.

\subsection{Question \& Answer}

At this stage a natural question is whether the order of \(\theta\) and \(\bar\theta\) carries any deep significance. Suppose we write \(c\,\bar\theta\theta\). Could we not just as well write \(c\,\theta\bar\theta\)?

Yes, but only at the price of a sign. Since both variables are odd,
\begin{equation}
\bar\theta\theta=-\theta\bar\theta.
\end{equation}
So we may rewrite
\begin{equation}
c\,\bar\theta\theta=-\,c\,\theta\bar\theta,
\end{equation}
or absorb the sign into the coefficient. If \(c\) is also odd, moving the variables past \(c\) produces further sign changes. What does not change is the parity. These manipulations may flip a sign, but they do not turn odd into even or even into odd. In the lecture's own register, this is still just bookkeeping for fermions. The deeper significance will come later.

\section{Ordinary Symmetry as the Comparison Case}

Only after the Grassmann prelude does the lecture return to the familiar language of continuous symmetry. Let \(G_i\) denote the infinitesimal generators of a continuous group. Their algebra closes under commutators:
\begin{equation}
[G_i,G_j]=i\,c_{ijk}G_k,
\end{equation}
where \(c_{ijk}\) are the structure constants. This closure is the infinitesimal version of the multiplication law of the group.

The generators of sensible physical symmetries are Hermitian. One may combine them into useful complex linear combinations without losing the underlying Hermitian interpretation. The usual angular momentum example is
\begin{equation}
L_\pm=L_x\pm iL_y.
\end{equation}
The point of recalling this is that later we will also use complex combinations of generators, but now the generators themselves will be odd.

For an actual symmetry of a physical system, the Hamiltonian must commute with the generators:
\begin{equation}
[H,G_i]=0.
\end{equation}
This is another way of saying that the symmetry does not change the energy. It also says that the corresponding charge is conserved.

Susskind emphasizes one more contrast before passing to supersymmetry. Ordinary measurable quantities are even. They are represented by ordinary numbers, or by operators whose eigenvalues are ordinary numbers. The ordinary symmetries of nature therefore live in an even algebra and are organized by commutators. Supersymmetry changes exactly this assumption: the generators themselves are odd. Once that happens, the natural algebraic operation is no longer the commutator but the anticommutator.

\section{A Minimal Supersymmetry Algebra from Oscillators}

To see what such a symmetry might look like, the lecture constructs the simplest possible toy model. We do not begin with a field theory in full space-time. We begin with one bosonic oscillator and one fermionic oscillator, both thought of as particles at rest. In that setting the bosonic creation and annihilation operators satisfy
\begin{equation}
[a_+,a_+]=[a_-,a_-]=0,
\qquad
[a_-,a_+]=1,
\end{equation}
and the fermionic operators satisfy
\begin{equation}
\{c_+,c_+\}=\{c_-,c_-\}=0,
\qquad
\{c_-,c_+\}=1.
\end{equation}
In particular,
\begin{equation}
c_+^2=c_-^2=0.
\end{equation}
The \(a\)'s are even, the \(c\)'s are odd, and every \(a\) commutes with every \(c\).

Now define the odd generators
\begin{equation}
Q=\sqrt{m}\,a_+c_-,
\qquad
\bar Q=Q^\dagger=\sqrt{m}\,a_-c_+.
\end{equation}
The lecture uses the notations \(Q^\dagger\) and \(\bar Q\) interchangeably; we will use \(\bar Q\) in the main flow. The action is direct and intuitive. \(Q\) removes a fermion and replaces it by a boson; \(\bar Q\) does the reverse. We are therefore looking at symmetry operations that exchange bosonic and fermionic states.

The crucial question is the analog of the Lie-algebra closure relation: what is \(\{Q,\bar Q\}\)? We write
\begin{align}
\{Q,\bar Q\}
&=
Q\bar Q+\bar Q Q \\
&=
m\bigl(a_+c_-a_-c_+ + a_-c_+a_+c_-\bigr).
\end{align}
Because the even \(a\)'s commute with the odd \(c\)'s, we may regroup:
\begin{equation}
\{Q,\bar Q\}
=
m\bigl(a_+a_-c_-c_+ + a_-a_+c_+c_-\bigr).
\end{equation}
Now use
\begin{equation}
c_-c_+=1-c_+c_-,
\qquad
a_-a_+=a_+a_-+1.
\end{equation}
Then
\begin{align}
\{Q,\bar Q\}
&=
m\Bigl[a_+a_-\bigl(1-c_+c_-\bigr)+\bigl(a_+a_-+1\bigr)c_+c_-\Bigr] \\
&=
m\bigl(a_+a_-+c_+c_-\bigr).
\end{align}
If we define
\begin{equation}
N_B=a_+a_-,
\qquad
N_F=c_+c_-,
\end{equation}
this becomes
\begin{equation}
\{Q,\bar Q\}=m(N_B+N_F).
\end{equation}

At this point the physical interpretation enters. In the toy model the particles are at rest, so the total energy is just the common mass times the total number of particles:
\begin{equation}
H=m(N_B+N_F).
\end{equation}
Thus the algebra closes as
\begin{equation}
\{Q,\bar Q\}=H.
\end{equation}

\begin{figure}
\centering
\includegraphics[width=0.62\textwidth]{lecture_05_figure_02.png}
\caption{Supersymmetry anticommutator yielding the Hamiltonian.}
\end{figure}

This is the first real climax of the lecture. We expected that anticommuting two odd generators might return some even element of the algebra. What we actually get is the Hamiltonian, the generator of time translations. The symmetry is therefore not merely an internal boson--fermion exchange. It is already entangled with space-time.

The remaining relations are immediate:
\begin{equation}
\{Q,Q\}=0,
\qquad
\{\bar Q,\bar Q\}=0,
\end{equation}
because \(Q^2\) contains \(c_-^2\) and \(\bar Q^2\) contains \(c_+^2\).

\subsection{Question \& Answer}

Why should \(Q\) commute with \(H\), and why do the boson and fermion masses have to be equal?

The lecture first gives the result and then pauses to justify it carefully. Let \(|E\rangle\) be an eigenstate of the Hamiltonian:
\begin{equation}
H|E\rangle=E|E\rangle.
\end{equation}
Assume now that the boson and fermion have the same mass. Acting with \(Q\) removes one particle and inserts another of exactly the same mass. So \(Q|E\rangle\) has the same energy \(E\). Therefore
\begin{align}
HQ|E\rangle &= E\,Q|E\rangle, \\
QH|E\rangle &= Q(E|E\rangle)=E\,Q|E\rangle.
\end{align}
Subtracting the two,
\begin{equation}
[Q,H]|E\rangle=0.
\end{equation}
If the energy eigenstates form a complete basis, it follows that
\begin{equation}
[Q,H]=0.
\end{equation}
The same reasoning applies to \(\bar Q\).

This is the point in the lecture where one hears the usual symmetry logic returning in a new form. A symmetry generator commutes with the Hamiltonian precisely because it does not change the energy. Here the new odd generator has that property only because boson and fermion masses have been matched. If the masses were different, the transformed state would not have the same energy, the commutator with \(H\) would fail to vanish, and the algebra would cease to close neatly. That is why the equal-mass condition is not an ornament. It is the condition under which the toy supersymmetry works at all.

\section{From the Hamiltonian to Superspace}

Once the Hamiltonian appears on the right-hand side of the supersymmetry algebra, the lecture takes one more backward step. If \(\{Q,\bar Q\}=H\), then we had better remind ourselves what \(H\) does. In this simple example there is no space, only time. But that is enough: if there is energy, there is time evolution.

For a state vector, the standard infinitesimal relation is
\begin{equation}
i\,\frac{\partial}{\partial t}|\psi\rangle = H|\psi\rangle.
\end{equation}
The board evidence at this stage shows only the time-derivative operator acting on the ket, not the full equation, so we keep the partial board expression nearby:
\begin{equation}
\frac{\partial}{\partial t}|\psi\rangle.
\end{equation}

\begin{figure}
\centering
\includegraphics[width=0.56\textwidth]{lecture_05_figure_03.png}
\caption{Time derivative acting on a state vector.}
\end{figure}

The same time evolution may be written in finite form:
\begin{equation}
e^{-iH\delta}|\psi(t)\rangle = |\psi(t+\delta)\rangle.
\end{equation}
The lecture now turns the same idea into a statement about a function. Suppose we consider
\begin{equation}
\psi(\theta,\bar\theta,t),
\end{equation}
a wavefunction or field-like object depending not only on time but also on Grassmann coordinates. An ordinary time translation is
\begin{equation}
t\to t+\delta,
\end{equation}
and to first order it induces
\begin{equation}
\psi(\theta,\bar\theta,t)\to
\psi(\theta,\bar\theta,t)+\delta\,\frac{\partial \psi(\theta,\bar\theta,t)}{\partial t}.
\end{equation}

\begin{figure}
\centering
\includegraphics[width=0.72\textwidth]{lecture_05_figure_04.png}
\caption{First-order time translation of the wavefunction.}
\end{figure}

This is still ordinary quantum mechanics. The radical move comes next. If space translations are generated by derivatives with respect to spatial coordinates, could there be an analogous symmetry associated with shifts of Grassmann coordinates? Susskind proposes that we consider
\begin{equation}
\theta\to\theta+\epsilon,
\qquad
t\to t+i\,\bar\theta\,\epsilon,
\end{equation}
where \(\epsilon\) is itself an odd Grassmann parameter.

This is a very strange transformation, and the lecture says so explicitly. But it is strange in a controlled way. Since \(\epsilon^2=0\), everything is automatically first order. If we shift \(\theta\) by \(\epsilon\), then
\begin{equation}
\delta\psi = \epsilon\,\partial_\theta\psi
\end{equation}
from the change in the Grassmann argument alone. The simultaneous change in \(t\) contributes another first-order term proportional to \(\bar\theta\,\partial_t\psi\). In a conventional normalization, one may summarize the total variation as
\begin{equation}
\delta_\epsilon\psi
=
\epsilon\Bigl(\partial_\theta-i\,\bar\theta\,\partial_t\Bigr)\psi.
\end{equation}
This suggests an odd differential operator
\begin{equation}
q=\partial_\theta-i\,\bar\theta\,\partial_t.
\end{equation}
Its conjugate partner is then written as
\begin{equation}
\bar q=\partial_{\bar\theta}+i\,\theta\,\partial_t.
\end{equation}

Here we should be careful, because the lecture itself is careful. Susskind says outright that he may have the sign wrong and that a factor of two may have slipped somewhere into the normalization. So the exact formulas above should be read as a standard reconstruction of the intended structure, not as a claim that every sign is visible on the board.

What is solid is the algebraic mechanism. When we take the anticommutator \(\{q,\bar q\}\), the derivative-derivative term vanishes, just as ordinary derivatives commute. The odd-coordinate--odd-coordinate term vanishes as well. The surviving terms are the mixed ones, and they use precisely the relations
\begin{equation}
\{\partial_\theta,\theta\}=1,
\qquad
\{\partial_{\bar\theta},\bar\theta\}=1.
\end{equation}
Thus the anticommutator is proportional to the time-translation generator:
\begin{equation}
\{q,\bar q\}\propto -\,i\,\partial_t.
\end{equation}
Up to the sign and factor-of-two uncertainty flagged in the lecture, this is the differential-operator version of the earlier relation \(\{Q,\bar Q\}=H\).

Now the earlier surprise becomes geometric. In the oscillator model the algebra closed on the Hamiltonian. In superspace the same fact is reinterpreted as follows: translations in the odd directions \(\theta\) and \(\bar\theta\), when anticommuted, produce an ordinary translation in time. Supersymmetry is therefore not merely a device for relabeling particles. It enlarges the geometry on which the theory lives.

\section{Summary}

The lecture unfolds in a very deliberate order. First we are taught how to manipulate odd variables: anticommuting Grassmann numbers, their truncated functions, their derivatives, and the basic identity
\begin{equation}
\{\partial_\theta,\theta\}=1.
\end{equation}
Only then are we reminded of the ordinary algebra of continuous symmetries and invited to change one decisive feature: let the generators themselves be odd.

In the minimal boson--fermion oscillator model this leads to
\begin{equation}
\{Q,\bar Q\}=H,
\end{equation}
with \(\{Q,Q\}=\{\bar Q,\bar Q\}=0\) and \([Q,H]=0\) when the boson and fermion masses are equal. The remarkable point is that the algebra closes only when the Hamiltonian is included. That is what ties supersymmetry to space-time.

The lecture then pushes this one step further. Since \(H\) generates time translations, the same algebra may be viewed geometrically. A shift in odd coordinates, together with an induced shift in time, leads to differential supercharges whose anticommutator reproduces time translation. In one sense this remains a bookkeeping device for bosons and fermions. In another sense it is a deep enlargement of the very notion of space-time.